Better forecasts can expand useful acceptance, but this alone does not specify how to spend the error budget. If a predictor's full-coverage risk is at most $r$, accepting all cases is a common optimum and mismatch disappears. The informative regime is therefore a predictor that improves while its full-coverage risk remains above the cap.

We compare a hierarchy-feature L1 forecaster with category scaling, its forward-time learned residual correction, and one fixed univariate Chronos-2 configuration \citep{ansari2025chronos2}. All are evaluated on the same consumed development rows. For this experiment only, we fit one error head per predictor plus a common demand head on days 1555--1610, calibrate on 1611--1645 and 1646--1673, and evaluate on 1800--1913. The same 600,000 rows, 220-tree specification, $r=0.64013$, and 0.35 row floor are used throughout. No old error head or threshold is reused. Chronos weights remain frozen; its selector is task-trained. This tests acceptance within each pipeline, not training-matched backbone superiority.

\input{accuracy_table.tex}
The residual learner lowers full-coverage WAPE by 1.76\% relative to category-scaled L1. Both its row-weighted and mixed policies increase \emph{both} coverages over the corresponding base policies. Yet switching utility within that improved predictor still gains 5.87 demand-coverage points while losing 3.02 row points. The paired item-bootstrap intervals are $[4.84,6.88]$ and $[-3.57,-2.43]$ points, respectively. The same signs occur for the base and Chronos predictors, and for all three pure-demand endpoints (Appendix Table~\ref{tab:accuracy-ci}). These are retrospective, fit-conditional comparisons.

\begin{figure}[t]
\centering\includegraphics[width=\linewidth]{figures/accuracy_coverage.pdf}
\caption{New matched-date comparison. Left: moving to a stronger predictor raises both coverages, while changing utility within a predictor trades them off. Circles, squares, and triangles denote $\lambda=1,0.25,0$. Lines connect evaluated settings, not Pareto frontiers. Right: the improved predictor's exclusive acceptance sets retain the budget mechanism.}
\label{fig:accuracy}
\end{figure}
Budget accounting connects the result to the theorem. For the improved predictor, the common set supplies 43,525.76 units of slack. The row-only set spends 26,046.82, accepting 135,756 rows with mean forecast 0.107, WAPE 0.9904, and only 2.55\% of total demand. The mixed-only set instead adds 72,762 rows carrying 8.42\% of demand. With the improved predictor, the row-coverage sacrifice is smaller than for the base (5.77 versus 3.02 points), while the different allocations remain. Hobbies and Household still miss the cap under every pooled policy; this experiment establishes no groupwise deployment guarantee.
